\section{Correctness findings}
\label{sec:correctness}

Every arm was checked on every call against a float64 reference of its own contract.
Table~\ref{tab:correctness} (Appendix~\ref{app:tables}) lists, per CPU and arm, the largest implementation error, the input-rounding effect and the top-1 agreement with the contract-S reference.
Apart from maxsim-cpu on the M3 Max (issue \#5 below), every exact arm with float32 inputs stayed within $4\times10^{-7}$ of its own reference on every call.
The checks found two new defects in other libraries and confirmed a third that was already reported. We reported the two new ones upstream in October 2026. The NumKong maintainer acknowledged the AMX defect, and another contributor opened a fix~\cite{numkong_pr409}; the maxsim-cpu report is still open.

\paragraph{maxsim-cpu, x86 build, long queries (issue \#1).}
For queries with more than 32 tokens, once a document has 64 or more tokens, the x86 wheel of maxsim-cpu~0.1.0 returned wrong scores, up to $10^{38}$ and infinity.
Another user had reported this behaviour in 2025~\cite{maxsimcpu_issue1}; our reference check found it independently, and we reproduced it on a second x86 CPU.
Its sums over 32-token query chunks are exact, so the benchmark calls it that way on x86; fewer than 1\% of agent calls have more than 32 query tokens.

\paragraph{maxsim-cpu, AArch64 build, zero padding (issue \#5).}
Its list interface pads shorter documents with zero vectors inside a call.
A zero vector scores 0, so for a query token whose best cosine is negative, the padded documents get 0 instead: a $\max(0,\cdot)$ clamp that depends on which other documents share the call~\cite{maxsimcpu_issue5}.
The ragged-length fixture catches it on AArch64.
On the M3 Max its largest error against its own contract was $4\times10^{-2}$ to $8\times10^{-2}$; its top document still agreed with the contract-S reference on every call, because S applies the same clamp on purpose.

\paragraph{NumKong, AMX path (issue \#406).}
On CPUs with AMX, NumKong~7.8.5's packed MaxSim kernel scans all 16 columns of a partially filled last tile.
When a document's token count modulo 64 is 49 to 63 and every cosine of a query token is negative, a padding column wins, and the result is wrong (NaN, 0, 1 or very large values)~\cite{numkong_issue406}.
The negative-best-match fixture found it on the Granite Rapids machine, and a standalone script reproduced it on a second machine of the same type.
The defect also occurred on the replayed data. On Intel, NumKong's float16 arms returned NaN scores on \nm{nan.intel.numkong_native_f16.ho-warm.agent-ho}\% of agent calls, \nm{nan.intel.numkong_native_f16.ho-warm.gcsn-ho}\% of GCSN calls and \nm{nan.intel.numkong_native_f16.ho-warm.beir}\% of BEIR calls; its float32 arms did so on up to \nm{nan.intel.numkong_native_f32.ho-warm.sweep}\% of sweep calls. No arm returned NaN on the CPUs without AMX. The defect's other wrong values (0, 1, very large) are not NaN and can hide inside NumKong's approximation error, so the Intel NumKong rows include an unknown number of further defective calls. Across runs, NumKong's Python arms on Intel also returned different scores for the same call (values up to $10^{74}$), while every other arm on every CPU returned bit-identical scores in all its runs. Their timings stand as measured and are marked ($\ast$) in the speed tables; their error and quality numbers do not describe NumKong's intended behaviour, so we take NumKong's quality from the AMD run.

\paragraph{Contract differences are not defects.}
NumKong's cosine of the received vectors (contract C) and PyLate's zero-padding rule (contract S-pre-pad0) are design choices, not bugs. maxsim-cpu's zero padding counts as a defect because the same call gives different scores on its x86 and AArch64 builds and depends on the other documents in the call.
Table~\ref{tab:contract-effects} measures them reference against reference.
Contract C keeps the top document of contract S on \nm{ce.agent-ho.S-vs-C.top1}\% of agent calls and \nm{ce.gcsn-ho.S-vs-C.top1}\% of GCSN calls, but on only \nm{ce.beir.S-vs-C.top1}\% of BEIR calls.
The clamp changes the top document on at most 0.1\% of calls; the epsilon placement changes no score by more than $10^{-15}$.
So almost all of the S--C difference comes from normalization.
Both models produce token vectors of unit length, so the stored norm $n_j$ is 1 (Section~\ref{sec:contracts}).
On the float vectors, S and C therefore differ only by the epsilon and the clamp.
On int4 tokens, C divides by the decoded norm, and S does not.
On development data, the decoded norm differed from 1 by up to about 2.5\%.
The two rules diverge most on BEIR, where candidate scores lie closest together.

\begin{table}[!htbp]
\centering\scriptsize
\caption{Contract effects, float64 reference against float64 reference (held-out sets; calls with at least two candidates). No implementation error enters these numbers.}
\label{tab:contract-effects}
\begin{tabular}{llrrrr}
\toprule
set & effect & calls & max $|\Delta|$ & top-1 agreement & top-10 overlap \\
\midrule
\inputtable{tables/contract_effects.tex}
\bottomrule
\end{tabular}
\end{table}
